% Generated by Listas/make_statements.py from pratica-lecture-03.md, item 9. Do not edit.
\textbf{Source.} Tadelis (2013), \emph{Game Theory: An Introduction}, Exercise 7.9 ``The Dean's Dilemma'', pp.~149--150 {[}PDF 168--169{]}. Reproduced verbatim. Answer in the source: none in the book (the preface, p.~xiii, says about half are solved on Princeton UP's website; not checked).



\begin{quote}

\textbf{The Dean's Dilemma:} A student stole the DVD player from the student lounge. The dean of students (player 1) suspects the student (player 2) and begins collecting evidence. However, evidence collection is a random process, and concrete evidence will be available to the dean only with probability \(\frac{1}{2}\). The student knows that the evidence-gathering process is under way but does not know whether the dean has collected evidence or not. The game proceeds as follows: The dean realizes whether or not he has evidence and can then choose his action, whether to accuse the student (\(A\)) or bounce the case (\(B\)) and forget it. Once accused, the student has two options: he can either confess (\(C\)) or deny (\(D\)). Payoffs are realized as follows: If the dean bounces the case then both players get 0. If the dean accuses the student and the student confesses, the dean gains 2 and the student loses 2. If the dean accuses the student and the student denies, then payoffs depend on the evidence: If the dean has no evidence then he loses face, which gives him a loss of 4, while the student gains glory, which gives him a payoff of 4. If, however, the dean has evidence then he is triumphant and gains 4, while the student is put on probation and loses 4.



\begin{enumerate}

\def\labelenumi{\alph{enumi}.}

\item

  Draw the game tree that represents the extensive form of this game.

\item

  Write down the matrix that represents the normal form of the extensive form you drew in (a).

\item

  Solve for the Nash equilibria of the game.

\end{enumerate}

\end{quote}

